\chapter{Febrile Seizure}

\section*{Definition}
A seizure occurring in children aged 6 months to 5 years, triggered by fever (temperature above 38	extdegree{}C) in the absence of intracranial infection, metabolic disturbance, or history of afebrile seizures. Simple febrile seizures last less than 15 minutes and are generalized; complex febrile seizures are focal, prolonged, or recurrent within 24 hours.

\section*{What Happens in Your Body}
Fever causes rapid temperature elevation that lowers the seizure threshold in the developing brain through immune signaling proteins-mediated effects (IL-1, IL-6) on neuronal excitability and GABAergic inhibition. Genetic susceptibility involving sodium channel and GABA receptor genes increases risk. The immature blood-brain barrier and developing myelination make young children particularly vulnerable. Simple febrile seizures do not cause brain damage.

\section*{Causes}
\begin{itemize}
  \item Upper respiratory tract infections (otitis media, pharyngitis, tonsillitis)
  \item Viral infections (influenza, adenovirus, RSV, HHV-6 roseola, enterovirus)
  \item Pneumonia or bronchiolitis
  \item Urinary tract infection (especially in young children)
  \item Gastroenteritis
  \item Bacterial infections (pneumonia, occult bacteremia)
  \item Post-immunization fever (MMR, DTaP, pneumococcal, influenza)
  \item Roseola infantum (HHV-6) — classic cause of febrile seizures
  \item Dental abscess or other focal infection
  \item Genetic predisposition (family history of febrile seizures)
\end{itemize}

\section*{When to See a Doctor}
See your doctor if:
\begin{itemize}
  \item Symptoms are severe or getting worse over time
  \item First febrile seizure always requires emergency evaluation to exclude meningitis, encephalitis, or electrolyte disturbance; seizure lasting more than 5 minutes, focal seizure, or multiple seizures within 24 hours
  \item You have other concerning symptoms such as fever, weight loss, or bleeding
\end{itemize}

\section*{Related Symptoms}
\begin{itemize}
  \item Fever
  \item Seizures
  \item Loss of consciousness
  \item Confusion
  \item Shaking
\end{itemize}
\textbf{Important:} If this symptom persists, your doctor will check for serious underlying causes.

\section*{Self-Care / Home Management}
\begin{itemize}
  \item During a seizure, place the child on their side in the recovery position and do not restrain them.
  \item Do not put anything in the child's mouth during a seizure.
  \item Remove nearby objects that could cause injury.
  \item Administer antipyretics (acetaminophen or ibuprofen) for fever control after the seizure.
  \item Call emergency services if the seizure lasts more than 5 minutes or is the child's first seizure.
\end{itemize}

\section*{How Your Doctor Will Evaluate You}
\begin{itemize}
  \item Your doctor will take a detailed history of the seizure, fever pattern, and recent illnesses.
  \item A full physical examination checks for the source of fever (ears, throat, chest, urine).
  \item Lumbar puncture may be performed to rule out meningitis or encephalitis if signs of CNS infection are present.
  \item Blood tests evaluate glucose, electrolytes, and complete blood count.
  \item EEG or brain imaging is not routinely indicated for simple febrile seizures but may be done for complex cases.
\end{itemize}

\section*{Treatment Options}
\begin{itemize}
  \item Antipyretics (acetaminophen, ibuprofen) to control fever and improve comfort.
  \item Treatment of the underlying infection (antibiotics for bacterial infections, supportive care for viral).
  \item Antiseizure medications are not routinely prescribed for simple febrile seizures.
  \item Rectal diazepam or intranasal midazolam may be prescribed for prolonged or recurrent seizures.
  \item Parental education and reassurance are key — simple febrile seizures do not cause brain damage.
\end{itemize}


\section*{References}
\begin{thebibliography}{99}
\bibitem{patel2015}
Patel N, Ram D, Swiderska N, Mewasingh LD. Febrile seizures. \textit{BMJ}. 2015;351:h4240.
\bibitem{millar2003}
Millar JS. Evaluation and treatment of the child with febrile seizure. \textit{Am Fam Physician}. 2006;73(10):1761--1768.
\bibitem{greenberg2014}
Greenberg SM, Rivkees SA, Friedman AH, Bhardwaj RD. \textit{Febrile Seizures}. In: \textit{Pediatric Neurology}. Springer; 2014.
\bibitem{american2011}
American Academy of Pediatrics, Subcommittee on Febrile Seizures. Neurodiagnostic evaluation of the child with a simple febrile seizure. \textit{Pediatrics}. 2011;127(2):389--394.
\bibitem{steinlin2005}
Steinlin M. Febrile seizures. \textit{Swiss Med Wkly}. 2005;135(19-20):282--289.
\bibitem{shinnar2002}
Shinnar S, Glauser TA. Febrile seizures. \textit{J Child Neurol}. 2002;17(suppl 1):S44--S52.
\end{thebibliography}
